\scanpage{72}
opération extérieure de $\operatorname{Gal}(\overline{\mathbb Q}/\mathbb Q)$ sur le premier mb de (38), et en particulier sur $\widehat\Pi_{0,3}$ lui-même...).

Les cas les plus intéressants sont les \deleted{deux} suivants :

I) $\nu_0=\nu_1=\nu_\infty=+\infty$ : on retrouve $\Pi_{0,3}^{\tau}$ et $\Pi_{0,3}^{\deleted{\unsure{\tau}}}$, classifiant les cartes triangulées pondérées (resp. orientées).

II) $\nu_0=\nu_\infty=+\infty$, $\nu_1=2$ : on retrouve les groupes $\Pi_{0,3}^{\tau}/\langle\rho_1^2\rangle$ et $\Pi_{0,3}/\langle\rho_1^2\rangle$, classifiant les cartes ordinaires.

III) $\nu_0=\infty$, $\nu_1=2$, $\nu_\infty=3$ : donnent les groupes $\Pi_{0,3}^{\tau}/\langle\rho_1^2,\rho_\infty^3\rangle$ et $\Pi_{0,3}/\langle\rho_1^2,\rho_\infty^3\rangle$, classifiant les cartes ordinaires \uline{triangulées} (\uline{pas} pondérées !), i.e. dont toutes les faces sont \deleted{des monogones \unsure{ou}} d'ordre divisant $3$ i.e. sont des monogones ou des triangles -- dans le cas non orienté, et dans le cas orienté. On trouve des isomorphismes remarquables (via (35)(36))

\margin{Ces isomorphismes sont valables pour tte situation d'un groupe d'opérateurs $G$ fidèle sur une surface topol. $X$, avec groupes d'isotropie finis, \unsure{en interprétant} $\Pi_1(X,G;P)$ mixte, en termes \deleted{\unsure{topologique}} d'un groupe fondamental ``avec signature'' $(\Sigma_0)$.}

\diagram{\begin{tikzcd}[ampersand replacement=\&,column sep=large,row sep=large]
\text{(40)}\ \raisebox{-1.1cm}{$\left\{\rule{0pt}{2.8cm}\right.$}
\&
\underbrace{\Pi_{0,3}^{\tau}/\langle\rho_1^2,\rho_\infty^3\rangle}_{\text{groupes cartographiques}}
\&
\underbrace{\Pi_1(U_{0,3},\unsure{\mathcal G_3|\mathfrak S_3}\times\{1,\tau\};P_+)}_{\Pi_{0,3}^{D_3,\tau}}
\arrow[l,"\sim"']
\\
\&
\unsure{=}\ \underbrace{\Pi_{0,3}/\langle\rho_1^2,\rho_\infty^3\rangle}_{\substack{\uline{\text{triangulés}},\ \unsure{\text{non or. resp.}}\\\text{orienté}}}
\arrow[r,leftrightarrow,"\sim"]
\&
\overbrace{\Pi_1(U_{0,3},\unsure{\mathcal G_3|\mathfrak S_3};P_+)}^{\Pi_{0,3}^{D_3}}
\end{tikzcd}}
\margin{$\{\rho_1,\rho_\infty\mid\rho_1^2=\rho_\infty^3=1\}$}
\margin{71}
